\scanpage{26}

\mspage{12}

\[
\left\{\begin{aligned}
\widetilde{SF}_M&\simeq\underbrace{\mathrm{Bsp}(M,e_M)}_{\{(x,y,z)\in B(M)\mid x\wedge y\in\{\pm e_M\}\}}/
\underbrace{\unsure{\mathbb{S}_3|\mathscr{S}_3|\mathfrak{S}_3}}_{\text{opérant librement}}\\
&\text{de cardinal}\quad 2N/6=N/3
\end{aligned}\right.\tag*{(60)}
\]

d'où

\[
\operatorname{Card}SF_M=\deleted{\unsure{\tfrac16}}1/2\ \operatorname{Card}\widetilde{SF}_M=N/6=N'/3\;,\tag*{(61)}
\]

en résumé on aura

\[
\left\{\begin{aligned}
s_0(SP_M)&=N'/\deleted{\unsure{2}}n\\
s_1(SP_M)&=N'/\deleted{\unsure{n}}2\\
s_2(SP_M)&=N'/3
\end{aligned}\right.\tag*{(62)}
\]
\margin{NB tous les sommets sont d'ordre $n$\\ toutes les faces sont d'ordre $3$}

% slanted
\margin{C'était évident d'avance, en vertu du th. et de l'analyse faite de la structure de $SP_M$ autour d'un sommet\ldots}

avec

\[
\left\{\begin{aligned}
N'&=\operatorname{card}\mathrm{SL}(2,\mathbb{Z}/n\mathbb{Z})'=\tfrac12\overbrace{\prod_\alpha p_\alpha(p_\alpha-1)(p_\alpha+1)p_\alpha^{3(\nu_\alpha-1)}}^N\\
&\text{où}\quad n=\prod_\alpha p_\alpha^{\nu_\alpha}\;,
\end{aligned}\right.\tag*{(63)}
\]

d'où

\[
\begin{aligned}
\chi(SP_M)&=N'\left[\frac1{\deleted{\unsure{2}}n}-\frac1{\deleted{\unsure{n}}2}+\frac13\right]=N'\left[\frac{\deleted{\unsure{5}}1}{\deleted{\unsure{6}}n}-\frac1{\deleted{\unsure{n}}6}\right]\\
&=2-2g
\end{aligned}\tag*{(64)}
\]

ou encore pour le genre de $SP_M$

\[
\deleted{g=1+\tfrac12N'[\unsure{\tfrac16}+\tfrac1n-\tfrac56]}\quad
g=1+\tfrac{N'}2\left(\tfrac16-\tfrac1n\right)\tag*{(65)}
\]

On a donc $g=0$, i.e. $g<1$, sss \deleted{$\unsure{1/n-5/6}<0$ i.e.} \deleted{$5/6-1/n<0$ i.e. sss} $\frac16-\frac1n<0$ i.e. \deleted{sss}

\[
\left\{\begin{aligned}
&\deleted{n\leq5\qquad\text{sss}\qquad g=0}\\
&\deleted{n=6\qquad\text{sss}\qquad g=1}\\
&\deleted{n\geq7\qquad\text{sss}\qquad g\geq2}
\end{aligned}\right.\tag*{\deleted{(\unsure{66|67})}\deleted{(\unsure{67|66})}}
\]

% archive sheet number: 25
